\section{Fixed-target moments and signed phase resources}
\label{v9:resources:section}

The next two sections keep the actual law, index and threshold fixed.
Throughout, $n\ge N=42$, $c=\CE$, and the attained fixed-target excess
$a$ is nonnegative. We assume the sharp predecessor
$C_{n-1}\le c$ and the contact and conservation
formulas~\eqref{v9:core:contact} and~\eqref{v9:core:conservation}.
The entry of Section~\ref{v9:entry:selected-route} supplies a complete
closed cover, an actual central gap $(-A,B)$, and bounds for the same
law's $\rho,\beta,\gamma,m,D$ and raw threshold $H$.
Here $D=(3+\sqrt{10})\rho-\gamma-3J$, with $J$ the actual weighted gap
sum. Write $\mathcal L=L_{\rm core}$ for the nonnegative leakage denoted
$L$ in the core conservation formula; this separates it from slope
bounds used below.
The physical noise remains $W=X-\E[X\mid B_0]$ and the count noise is
$U=W/\ell$. Below, the different symbol $\Omega$ denotes a signed
Fourier moment coefficient.

\begin{condition}[Panel input and finite contraction]\label{v9:resources:input}
Use the nonempty probability panels indexed 24 through 53 in the selected
entry, with closed endpoints $[i/100,(i+1)/100]$. On each panel $I$ retain
its actual-law bounds, including $nD\le ND_I$ and an initial noise ceiling
$\theta\le\Theta_I$. Every law allowed by the entry must belong to at
least one such panel. The support radius $R=157/50$ is the largest
already acquired radius of the surviving fourth-pass rho rows.
All arithmetic in the ordered iteration below is outward: lower roots
are rounded down, upper roots up, and only a strictly negative residual
excludes a panel. Perform eight rounds; check every denominator is
positive before division and every asserted root has nonnegative input.
\end{condition}

The numerical input is a finite obligation, not a replacement for its
entry hypotheses. The formulas below explain the implication of each
recorded inequality and permit a check of each entire panel.

Write \(p=\Pp(X>0)\), \(q=1-p\), \(m=\E X_+=\E X_->0\).
The conditional means are \(\bar B=m/p\), \(\bar A=m/q\), and
\[
 \theta=1-\frac{m^2}{pq},\qquad
 \ell=\frac m{pq}=\sqrt{\frac{1-\theta}{pq}},\qquad
 d=1-\frac{A+B}{\ell},\qquad
 \lambda=\frac{n\theta}{\ell^2}.
\]
Thus \(d\) is a count-coordinate inward displacement, not \(D\).
Let \(V=\bar B\) on the positive sign and \(V=-\bar A\) on the negative sign;
write \(W=X-V\). Define the two measures on nonnegative magnitudes by
\(\int f(x)d\mu_\pm(x)=\E[W^2\mathbf1_{\{\pm X>0\}}f(|X|)]\).
Their masses sum to \(\theta\); the central gap ensures there is no atom at zero.
Also \(A\le\bar A\), \(B\le\bar B\).

\subsection{A common defect budget}\label{v9:resources:moment-cost}
We use the actual-law outer-gap inequality
\[
 D\ge(\sigma-1)\E X_+^3+(\sigma+1)\E X_-^3
       -6mAB+6BR_++6AR_-,
\]
\[
 \sigma=\sqrt{10},\qquad R_+=\E X_+^2-mB,\qquad R_-=\E X_-^2-mA.
\]
Here is the gap argument behind this input. On the positive side, summing
over outer support gaps and comparing each squared gap length with
$2\int(t-B)\,dt$ gives
\[
 J_o^+\le2\int_B^\infty(t-B)\tau(t)\,dt
 =\E[X_+(X_+-B)^2]=T_+-2BR_+,
 \quad T_+=\E X_+^3-mB^2.
\]
The negative side gives $J_o^-\le T_--2AR_-$ by the same magnitude
calculation. Adding the central term $m(A+B)^2$ and substituting in
$D$ proves the displayed inequality. Endpoint atoms are included in
the tail integrals; changing their values at integration endpoints does
not change an integral. Substituting
\(\E X_+^3=m\bar B^2+\int(x+2\bar B)d\mu_+\) and
\(R_+=\mu_+(\R)+m(\bar B-B)\), and their negative-sign counterparts, gives
\begin{equation}\label{v9:resources:moment-budget}
 D\ge \mathfrak b_0+E_0+\int k_+d\mu_++\int k_-d\mu_-,
\end{equation}
where, with \(p_*=(4-\sqrt{10})/2\),
\[
\begin{split}
 \mathfrak b_0&=\frac{(6+2\sigma)(1-\theta)^{3/2}}{\sqrt{pq}}(p-p_*)^2,\\
 E_0&=6m\{(A+B)(\bar A-A+\bar B-B)+(\bar A-A)(\bar B-B)\},\\
 k_+(x)&=(\sigma-1)(x+2\bar B)+6B,\qquad
 k_-(x)=(\sigma+1)(x+2\bar A)+6A .
\end{split}
\]
The deterministic identity is
\((\sigma+1)p^2+(\sigma-1)q^2-6pq=(6+2\sigma)(p-p_*)^2\).
Every displayed cost is nonnegative.

\subsection{Ordered contractions on closed probability panels}\label{v9:resources:iteration}
For $I=[l,u]$, set \(v_-=\min_{p\in I}pq\), \(v_+=\max_{p\in I}pq\).
The maximum includes \(1/4\) if \(I\) crosses \(1/2\).
Use \(\sigma_l=3.162277660<\sigma<3.162277661=\sigma_u\).
Initially \(D\le D_I<\sigma_l-1\), and the gap bound
\(J\le\rho+2mAB\) gives
\[
 AB\ge z:=(\sigma_l-1-D_I)/3>0.
\]
Indeed \(\rho\ge1,\gamma\le\rho,m\le1/2\) imply this inequality.
With \(\bar A_u=\sqrt{u/(1-u)}\), \(\bar B_u=\sqrt{(1-l)/l}\), obtain
\[
 A\ge A_l=z/\bar B_u,\quad B\ge B_l=z/\bar A_u,\quad
 (A+B)/\ell\ge g=2\sqrt{zv_-},\quad d\le1-g.
\]
For a currently proved \(\theta\le\Theta\), lower-bound the binary charge by
\[
 b_I=\frac{(6+2\sigma_l)(1-\Theta)^{3/2}}{\sqrt{v_+}}
        \operatorname{dist}(I,[ (4-\sigma_u)/2,(4-\sigma_l)/2])^2.
\]
All square roots are directed outward. If \(r=D_I-b_I<0\), the entire
closed panel is impossible. Otherwise set
\[
 \bar A_l=\sqrt{(1-\Theta)l/(1-l)},\quad
 \bar B_l=\sqrt{(1-\Theta)(1-u)/u},
\]
\[
 K=\min\{(\sigma_l-1)(B_l+2\bar B_l)+6B_l,\
          (\sigma_l+1)(A_l+2\bar A_l)+6A_l\}>0.
\]
Equation~\eqref{v9:resources:moment-budget} gives \(\theta\le r/K\).
Intersect this bound with \(\Theta\), then recompute \(b_I,r\).
Put \(m_l=\sqrt{(1-\Theta)v_-}\) and
\(\ell_l=\sqrt{(1-\Theta)/v_+}\). Since
\(E_0\ge6m\ell^2gd\), we may update
\[
 d\le\min\{d_{\rm old},r/(6m_l\ell_l^2g)\},\quad
 g\ge1-d,
\]
\[
 A_l\gets\max\{A_l,\ell_l(l-d)_+\},\quad
 B_l\gets\max\{B_l,\ell_l(1-u-d)_+\}.
\]
Only previously proved bounds enter each update. Zero residual forces
zero noise and displacement; no division by that residual is needed.
The implementation performs eight rounds, rounding upper bounds upward and
lower bounds downward to \(10^{-10}\).

\subsection{Signed moments and a weighted noise bound}\label{v9:resources:exports}
For every orientation, including \(p>1/2\),
\begin{equation}\label{v9:resources:beta}
 1-2u-2(1-u)\Theta\le\beta\le1-2l+2l\Theta.
\end{equation}
This follows from
\(\beta=(1-\theta)(1-2p)+\mu_+(\R)-\mu_-(\R)\).
The formerly convenient upper replacement \(1-2l+\Theta\) requires
\(l\le1/2\) and must not be continued past that point.

For a coherent third-moment cap, put
\[
 L=R+2\max(\bar A_u,\bar B_u),\quad
 k=\min\{\sigma_l-1+6B_l/(R+2\bar B_u),\
          \sigma_l+1+6A_l/(R+2\bar A_u)\}.
\]
The exact identity
\[\rho=\rho_B(p)(1-\theta)^{3/2}+\int(x+2\bar B)d\mu_++\int(x+2\bar A)d\mu_-.\]
gives a noise charge at most \(\min(L\theta,r/k)\).
Let \(r_B=(1-2v_-)/\sqrt{v_-}\) and
\(s\le(3/2)\sqrt{1-\Theta}\). Integrating the derivative proves
\((1-\theta)^{3/2}\le1-s\theta\) for \(0\le\theta\le\Theta\).
The maximum of this affine, piecewise expression is therefore bounded by
\[
 \rho\le r_B+(L-sr_B)_+\min\{\Theta,r/(kL)\}.
\]
Intersect with the prior rho cap.

Finally \eqref{v9:resources:moment-budget} implies the joint inequality
\[
 D\ge b_I+6m(A_l+B_l)\ell d+K\theta.
\]
Multiplying by the actual \(n\), using \(nD\le ND_I\),
\(m\ell=1-\theta\), and \(n\theta=(1-\theta)\lambda/(pq)\),
proves, when \(r>0\),
\[
 \frac{d}{r/[6(1-\Theta)(A_l+B_l)]}
 +\frac{\lambda}{Nv_+r/[K(1-\Theta)]}\le1.
\]
This retains the weighted resource; it is not obtained by multiplying an
arbitrary unweighted noise cap by \(N\).

\begin{theorem}[Conditional all-panel contraction]\label{v9:resources:all-panels}
Under Condition~\ref{v9:resources:input}, the exact eight-round
calculation on the selected inputs excludes panels 24 and 53.
On all 28 remaining panels,
\[
 \frac14\le p\le\frac{53}{100},\qquad
 \theta\le0.0594425514,\qquad d\le0.0719539227,
\]
\[
 \rho\le1.4547288545,\qquad \lambda\le0.6314459997,
\]
\[
 1.957063175\le\ell\le2.3094010768.
\]
All endpoints and both signs are included.
\end{theorem}
\begin{proof}
The previous inequalities justify each recorded rational update.
The input is the full accepted closed cover, and only a strictly negative
residual excludes a panel. All 28 retained outputs and every update are in
the exact log; the stated constants are their outward maxima or minima.
No sampled parameter or favorable moment optimizer is used.
\end{proof}

\subsection{Acquiring the incoming strip}\label{v9:resources:strip}
The retained fixed-target contact theorem, with sharp predecessor and
nonnegative excess, bounds a support contact by \(H_n(H,\pm t)+cQ(t)\),
where
\(Q(t)=-t^3+3rt^2/2+3bt+(8/875)r\).
Here \(r\) is a rho upper bound, and \(b=\beta_u\) on the positive ray,
\(b=-\beta_l\) on the negative ray.
For \(t>\sqrt3\), set
\[
 A_t=(t^3-3t)/6,\quad H_4=t^4-6t^2+3,\quad S=81/80,
\quad E_t=t^5/4100+3t/820+17(3t^2-2)/78720.
\]
The heat expansion used in the preceding entry proof remains valid here:
its fifth-derivative error has coefficient \(3t^2-2\), since \(t^2-1>0\).
Using \(\|\phi'''\|_\infty<2\phi(0)\), \(\phi(0)<2/5\) and
\(\sqrt{n-1}>6\), without any sign assumption on \(H_4\), yields
\[
 H_n(H,\pm t)\le \tfrac25S^3(A_t+|H_4|/72)+S^3E_t.
\]
Thus a negative boundary value of this expression plus
\((819/2000)Q(t)\) excludes that boundary, provided \(Q(t)<0\).
The exact derivative estimate is
\(|\partial_tH_n|<\tfrac25(t^2/2+1/2+8/875)\).
Check \(Q,Q',Q''<0\) and
\[
 L=\tfrac25(t^2/2+1/2+8/875)+(819/2000)Q',\quad L',L''<0.
\]
Their higher derivatives are negative constants, so both rays are excluded
beyond the boundary, for every raw threshold \(H\).
The complete Gaussian expansion and its domain are proved in
Lemma~\ref{v9:confinement:heat}; the finite boundary and continuation
inequalities remain explicit certificate conditions.

Choose the two boundaries
\(t_-=\ell_l(l+3/5)\) and \(t_+=\ell_l(1-u+3/5)\).
All satisfy \(t_\pm>\sqrt3\). Both pass on panels 40--52, proving
\(Y=X/\ell+p\in(-3/5,8/5)\) there.
On panels 25--39 the negative ray passes but the positive boundary
upper estimate is nonnegative. This is failure of that sufficient estimate,
not a counterexample to the strip theorem. No strip claim is made on those
fifteen panels by this test. On panels 40--52 the inward bound already
gives $U_0\le d<0.072$ and $U_1\ge-d>-0.072$; together with the
outer strip this proves $|U|<3/5$. This acquisition precedes the
half-amplitude argument in Section~\ref{v9:confinement:section}.

For later use, the panel $[0.41,0.42]$ has the following convenient
strict outward bounds. They are weaker than its eight-round output and
its retained entry phase bound:
\begin{equation}\label{v9:resources:central-entry}
\begin{gathered}
 \theta<0.054,\quad d<0.069,\quad\rho<1.22,\quad
 0.09736\le\beta\le0.234,\\
 1.9706<\ell<2.034,\qquad\lambda<0.6,\qquad |H|<1.814,
 \qquad |U|<0.6.
\end{gathered}
\end{equation}
For example, its recorded noise and displacement bounds are
$\theta\le0.053026061$, $d\le0.0672060498$,
$\lambda\le0.5728987925$ and $\rho\le1.2049635459$.
These numerical consequences still require the complete selected entry
and the finite contraction checks of Condition~\ref{v9:resources:input}.

\subsection{Signed Fourier terms from one actual moment cell}
\label{v9:resources:signed}
Write $M_j=\E|X|^j$ when an absolute moment is intended, and define
\[
 K_1=\E[X|X|^3]-3\beta-\tfrac32\rho\gamma,
 \qquad \Omega=J_1+(M_4-3)/6,
 \qquad J_1=\sum_{(a,b)}\tau_{a,b}\int_a^b(t-a)t\,dt.
\]
The moment coefficient $\Omega$ is unrelated to the physical noise $W$.
The complete Fourier estimate in the ordered entry supplies, on each
retained cell,
\begin{equation}\label{v9:resources:raw}
 \mathcal R_c\le K_F+cA_F+(cK_1+\varphi\Omega)H+s_F|H|
             -(c\rho+\varphi\gamma/2)H^2,
 \qquad\varphi=\phi(H/\sqrt n).
\end{equation}
Here $K_F,s_F\ge0$ are complete Fourier error bounds and $A_F$ bounds
the constant moment/gap term. This inequality retains all frequency
windows and their gates from Section~\ref{v9:entry:ordered-update};
the following algebra does not replace those analytic or finite inputs.
The cell contains ranges for $\rho,\gamma,K_1,\Omega,\varphi$ belonging
to this same law. Write their lower moment endpoints as $r_l,g_l$.

There are two ordered raw ceilings. The preceding pass supplies the
positive density floor $\varphi_l$ used in \eqref{v9:resources:raw}.
After bounding that polynomial one obtains a new ceiling
$\mathcal R_c\le T<Nc_-r_l$, where $c_-\le c\le c_+$.
Only now set
\[
 \eta=\frac{c_-r_l-T/N}{2}>0,\qquad
 \mathcal U=(n+H^2)a+n\phi(0)D/6+n^{3/2}\mathcal L.
\]
All three costs in $\mathcal U$ are nonnegative. The phase-price formula
of Corollary~\ref{v9:core:phase-price} gives
\begin{equation}\label{v9:resources:price}
 \mathcal U\le\mathcal R_c-\eta H^2.
\end{equation}
To see the ordering directly, put $w=H^2/(2n)$ and
$M=c\rho-\phi(0)D/6$. Conservation gives
$e^{-w}M\ge c\rho-T/n>0$, and therefore
\[
 \mathcal R_c-\mathcal U=n(1-e^{-w})M
 =n(e^w-1)e^{-w}M
 \ge\tfrac12H^2(c\rho-T/n)\ge\eta H^2.
\]
The density floor is not inferred from this new bound and then fed back
into the polynomial that produced it.

For each $t\in\{c_-,c_+\}$ define
\begin{align}
 B_t&=K_F+tA_F,&
 Q_t&=tr_l+\min_{v\in[\varphi_l,2/5]}vg_l/2+\eta,
       \label{v9:resources:quadratics}\\
 F_{t,+}&=\max(tK_1+v\Omega)+s_F,&
 F_{t,-}&=-\min(tK_1+v\Omega)+s_F.\notag
\end{align}
The extrema in the second line range over the recorded rectangular
intervals for $K_1,v,\Omega$. The expression is affine in each variable
separately, so its extrema occur at the eight corners. The separate
gate $Q_t>0$ must hold for both target endpoints.

\begin{proposition}[One-cell signed resource]\label{v9:resources:cell}
For $h=|H|$ and the actual sign $\varepsilon$ of $H$,
\begin{equation}\label{v9:resources:envelope}
 0\le\mathcal U\le
 \min\left\{T,\max_{t\in\{c_-,c_+\}}
                 (B_t+F_{t,\varepsilon}h-Q_th^2)\right\}.
\end{equation}
At $H=0$ both signed branches give the same value.
\end{proposition}
\begin{proof}
For fixed $h\ge0$, the maximum or negative minimum defining the signed
linear coefficient is a maximum of affine functions of $t$ and hence
convex. The remaining terms in the polynomial are affine in $t$.
Its maximum on $[c_-,c_+]$ is bounded by the larger endpoint value.
Apply \eqref{v9:resources:raw} and \eqref{v9:resources:price}; independently,
$\mathcal U\le\mathcal R_c\le T$. This proves the minimum with $T$.
In particular a negative signed slope remains negative. Replacing it by
zero is unnecessary and discards information.
\end{proof}

\subsection{Closed-region maxima and necessary noise filters}
\label{v9:resources:regions}
On $p\in[l,u]\subset(0,1)$ and $0\le\lambda\le\Lambda$, the exact
variance identities give
\[
 \theta=\frac{\lambda}{nv+\lambda},\qquad
 m^2=\frac{v^2}{v+\lambda/n},\qquad v=p(1-p).
\]
Thus with the actual minimum and maximum $v_-,v_+$ over the interval,
\begin{equation}\label{v9:resources:noise-filter}
\begin{gathered}
 \theta\le\Theta_\Lambda:=\frac{\Lambda}{Nv_-+\Lambda},\qquad
 \sqrt{\frac{v_-^2}{v_-+\Lambda/N}}\le m\le\sqrt{v_+},\\
 1-2u-2(1-l)\Theta_\Lambda\le\beta
       \le1-2l+2u\Theta_\Lambda.
\end{gathered}
\end{equation}
Indeed $v^2/(v+s)$ increases with $v>0$ and decreases with $s\ge0$.
The beta bounds follow by enlarging each coefficient in the exact signed
second-moment inequalities outward. This deliberately conservative filter
is valid on either side of $1/2$. It also includes the limit $n\to\infty$;
no finite upper bound on $n$ is needed.

Consider a closed rectangle $h\in[h_0,h_1]$,
$\lambda\in[\lambda_0,\lambda_1]$ on the probability panel. Retain
every moment cell whose $p,m,\beta$ intervals meet the necessary ranges
in \eqref{v9:resources:noise-filter}, with $\Lambda=\lambda_1$, and
whose previously acquired phase ceiling is at least $h_0$.
Only disjoint ranges can be discarded. In each retained cell clip $h_1$
by that phase ceiling. A concave quadratic $B+Fh-Qh^2$, $Q>0$,
has its maximum at the vertex clipped to this closed interval:
\[
 h_* =\max\{h_0,\min(h_1,F/(2Q))\}.
\]
Take the maximum of the two target-endpoint maxima, intersect with that
cell's $T$, and finally take the maximum over every retained cell.
Denote the result by $S_\varepsilon$. Then
\begin{equation}\label{v9:resources:rectangle}
 \mathcal U\le S_\varepsilon
\end{equation}
on the whole closed rectangle. Shared endpoints remain in every
intersecting cell; both phase signs are retained at zero. An empty list
excludes a rectangle only relative to the complete ancestor cover.
A negative $S_\varepsilon$ excludes it because $\mathcal U\ge0$.

\begin{condition}[Complete signed resource export]\label{v9:resources:phase-condition}
The selected fourth pass covers all 2,048 moment boxes and retains every
one of its 682 active cells. Both endpoint quadratics are reconstructed
on each active cell, giving 1,364 quadratics with their preceding density
floors and current positive phase prices. Rejected cells retain their
valid exclusion reasons; no missing or unclassified box is interpreted
as impossible. Every resource rectangle uses all incident active cells
and the complete signed maxima described above.
\end{condition}

This finite condition is checked by the exact replay and separate
coefficient reconstruction accompanying the selected entry. Its role is
to certify the actual-law cover and the rational bounds used here.
It does not prove contact, conservation or a Fourier remainder theorem.
Retaining the signed quadratics requires no new frequency integration;
it preserves information already contained in that last signed pass.
